\chapter{Pneumocystis jirovecii Pneumonia}
\index{Pneumocystis jirovecii Pneumonia}
\label{sec:Pneumocystis jirovecii Pneumonia}

\section{Overview}
Pneumocystis jirovecii pneumonia (PCP) is an opportunistic fungal infection of the lungs caused by the atypical yeast-like fungus Pneumocystis jirovecii, previously classified as a protozoan and formerly named Pneumocystis carinii. It remains a defining opportunistic infection in patients with advanced HIV/AIDS (CD4 count below 200 cells/$\mu$L) and also occurs in other immunocompromised populations including organ transplant recipients, patients on chronic corticosteroids, and those receiving immunosuppressive chemotherapy. Despite effective prophylaxis, PCP continues to be a significant cause of morbidity and mortality in immunocompromised patients worldwide.
\section{Classification}

\begin{description}[font=\normalfont\bfseries,leftmargin=3em,labelwidth=2.5em]
  \item[Cause] Infectious (fungal)
  \item[Organ System] Lungs/Respiratory
  \item[Pathophysiology] Attachment of Pneumocystis organisms to Type I alveolar epithelial cells causing foamy intra-alveolar exudate, alveolar-capillary barrier disruption, interstitial inflammation, and impaired gas exchange
  \item[Duration] Acute to Subacute
  \item[Onset] Gradual (over days to weeks)
  \item[Mode of Transmission] Airborne (inhalation of fungal elements)
  \item[Clinical Course] Progressive if untreated; respiratory failure develops in advanced cases
  \item[Severity] Moderate to Severe
  \item[Extent of Disease] Localized to lungs (can disseminate in severe immunosuppression)
  \item[Age Group] Immunocompromised individuals of all ages
  \item[Epidemiology] Leading AIDS-defining illness worldwide; incidence decreased with antiretroviral therapy and prophylaxis but remains significant in undiagnosed HIV and non-HIV immunosuppressed patients
  \item[ICD-11 Code] CA42.0
\end{description}

\section{Causes}
Pneumocystis jirovecii is an atypical fungus transmitted via airborne route that causes infection almost exclusively in immunocompromised hosts. In HIV-associated PCP, disease typically occurs when the CD4 count falls below 200 cells/$\mu$L. Non-HIV risk factors include prolonged high-dose corticosteroid use, solid organ or hematopoietic stem cell transplantation, hematologic malignancies, chemotherapy (especially rituximab and purine analogs), and primary immunodeficiency disorders such as severe combined immunodeficiency (SCID).

\section{Risk Factors}

\begin{itemize}[noitemsep]
  \item HIV/AIDS with CD4 count less than 200 cells/$\mu$L
  \item Chronic corticosteroid therapy (equivalent to prednisone $\ge$20 mg/day for $\ge$2 weeks)
  \item Solid organ transplantation (kidney, liver, lung, heart) with immunosuppression
  \item Hematologic malignancies (leukemia, lymphoma) and chemotherapy
  \item Rituximab, fludarabine, and other lymphocyte-depleting agents
  \item Primary immunodeficiency disorders (SCID, hyper-IgM syndrome)
  \item Inflammatory conditions requiring prolonged immunosuppression (Wegener granulomatosis, inflammatory bowel disease)
\end{itemize}

\section{Signs and Symptoms}

\subsection{Early symptoms}
\begin{itemize}[noitemsep]
  \item Early manifestations of pneumocystis jirovecii pneumonia may be subtle and nonspecific
\end{itemize}

\subsection{Common symptoms}
\begin{itemize}[noitemsep]
  \item \subsection{Common symptoms}
  \item \textbf{Common symptoms:}
  \item Subacute onset of progressive dyspnea on exertion
  \item Non-productive cough (dry cough)
  \item Fever with night sweats
  \item Bilateral diffuse pulmonary infiltrates on chest imaging
  \item Hypoxemia with increased alveolar-arterial oxygen gradient
  \item Tachypnea and labored breathing
  \item Fatigue and weight loss
  \item Chest tightness or discomfort
  \item Pain or discomfort in the affected area
  \item Difficulty performing daily activities
\end{itemize}

\subsection{Less common symptoms}
\begin{itemize}[noitemsep]
  \item Atypical or less common presentations of pneumocystis jirovecii pneumonia may occur
\end{itemize}

\subsection{Emergency warning signs}
\begin{itemize}[noitemsep]
  \item Severe or worsening symptoms of pneumocystis jirovecii pneumonia require immediate medical evaluation
\end{itemize}
\section{Progression and Stages}
PCP typically develops subacutely over 1--3 weeks in HIV-positive patients with progressive dyspnea, fever, and dry cough, while a more fulminant course over days is seen in non-HIV immunocompromised patients. Untreated infection progresses to severe hypoxemic respiratory failure requiring mechanical ventilation. With treatment, clinical improvement is usually evident within 4--8 days, though radiographic resolution may take weeks. Relapse can occur if immunosuppression persists without prophylaxis.

\section{Complications}

\begin{itemize}[noitemsep]
  \item Severe hypoxemic respiratory failure requiring intubation and mechanical ventilation
  \item Pneumothorax (spontaneous or from barotrauma)
  \item Acute respiratory distress syndrome (ARDS) in severe cases
  \item Disseminated infection (lymph nodes, liver, spleen, bone marrow) in profound immunosuppression
  \item Mortality of 10--20\% in HIV-associated PCP and 30--50\% in non-HIV immunocompromised patients
  \item Immune reconstitution inflammatory syndrome (IRIS) after starting antiretroviral therapy
  \item Reduced quality of life
  \item Emotional and psychological effects
\end{itemize}

\section{Diagnosis}

\begin{itemize}[noitemsep]
  \item Chest radiograph or CT showing bilateral interstitial or alveolar infiltrates (ground-glass opacification)
  \item Elevated serum lactate dehydrogenase (LDH) as a suggestive but non-specific marker
  \item Induced sputum or bronchoalveolar lavage (BAL) with Gomori methenamine silver stain, immunofluorescence, or PCR
  \item Bronchoscopy with BAL is the diagnostic gold standard
  \item Beta-D-glucan assay in serum (high sensitivity but not specific)
  \item Pulse oximetry and arterial blood gas to assess oxygenation
  \item HIV testing in newly diagnosed patients
  \item Imaging tests to evaluate extent of lung involvement
\end{itemize}

\section{Prevention}

\begin{itemize}[noitemsep]
  \item Primary prophylaxis with trimethoprim-sulfamethoxazole (TMP-SMX) for HIV patients with CD4 $<$200 cells/$\mu$L and other high-risk immunocompromised patients
  \item Alternative prophylactic regimens: dapsone, atovaquone, or aerosolized pentamidine
  \item Antiretroviral therapy (ART) for HIV patients to maintain CD4 count above 200
  \item Consider discontinuation of primary prophylaxis when CD4 count remains $>$200 for $\ge$3 months on ART
  \item Environmental precautions (avoiding crowded settings) for profoundly immunocompromised patients
  \item Go for regular check-ups
  \item Follow your doctor's screening recommendations
\end{itemize}

\section{Treatment Options}

\begin{itemize}[noitemsep]
  \item First-line therapy: trimethoprim-sulfamethoxazole (TMP-SMX) for 21 days
  \item Adjunctive corticosteroids for moderate-to-severe PCP with PaO2 $<$70 mmHg or alveolar-arterial gradient $>$35 mmHg
  \item Second-line therapy: atovaquone suspension or clindamycin plus primaquine
  \item Pentamidine (intravenous) for severe disease with contraindications to TMP-SMX
  \item Dapsone plus trimethoprim as alternative outpatient therapy
  \item Initiation or optimization of antiretroviral therapy in HIV patients
  \item Supportive care with supplemental oxygen and respiratory support as needed
  \item Lifestyle changes and self-care
  \item Surgery when needed (chest tube for pneumothorax)
\end{itemize}

\section{Living with It}

\begin{itemize}[noitemsep]
  \item Follow your treatment plan as prescribed
  \item Keep up with regular appointments
  \item Adjust your daily habits as needed
  \item Reach out to family, friends, or support groups
  \item Keep learning about your condition
\end{itemize}

\section{Outlook}
The prognosis for PCP depends on the underlying immune status with mortality of 10--20\% in HIV-associated cases and 30--50\% in non-HIV immunocompromised patients. With appropriate treatment and, in HIV patients, effective antiretroviral therapy, most patients recover from the acute infection and can resume normal activities with appropriate ongoing prophylaxis.
